\paragraph{The original ordinary modules at auxiliary levels.}
Fix the original definite unitary group over $F=L^+$, its Hida
setting, the coefficient representation $\tau$, and the continuous
weight character $\theta$. Let $(A,\mathfrak m)$ be the original
complete Noetherian coefficient algebra with finite residue field of
characteristic $\ell$. At the split places above $\ell$, retain the
source groups
\[
 I(c)=NT\overline N(c),\qquad
 \overline N(c)=\overline N\cap K(c+1),\qquad
 T_c=T\cap K(c+1),
\]
and its unnormalized operator
$U=[NzN]$, where
$z=\operatorname{diag}(\varpi^{n-1},\ldots,\varpi,1)$.
For an original tame level $K$, let $S(K,A)$ be the original
$A^m$-valued automorphic functions: left rational invariance, the
$\tau$-twisted right $K$ action, right $N$ invariance,
$f(gt)=\theta(t)^{-1}f(g)$, and uniform $\mathfrak m$-adic continuity
under the groups $\overline N(c)$. Write $F^aS$ for functions whose
values lie in $(\mathfrak m^a)^m$. Define
\[
 M(K,A)=\{f\in S(K,A):
       \ \forall a\ \exists r\geq1,\ U^rf-f\in F^aS\}.
 \tag{OD1}
\]

\begin{proposition}[Integral ordinary modules and level trace]
\label{prop:selmer-ordinary-diamond}
For every original compact auxiliary level $K$, the module $M(K,A)$
is finite free over $A$, and
$M(K,A)\cap F^aS=\mathfrak m^aM(K,A)$.
If $K'\triangleleft K$ has finite quotient $\Delta$, then
$M(K',A)$ is projective over $A[\Delta]$ and the original
unnormalized trace satisfies
\[
 \operatorname{Tr}_{K'}^K:M(K',A)\twoheadrightarrow M(K,A),
 \qquad
 \ker\operatorname{Tr}_{K'}^K
   =\sum_{\delta\in\Delta}(\delta-1)M(K',A).
 \tag{OD2}
\]
The assertions persist after the prescribed residual Hecke
localization.
\end{proposition}

\begin{proof}
Borel's finite class-number theorem gives finite double cosets at
every compact open level \cite{borel-ordinary-diamond}.
Their stabilizers are trivial for the actual Hida setting. Indeed,
the definite real group is compact, so a rational stabilizer is a
finite group. At the fixed place $v_0$ it is conjugate into
$1+pM_n(\mathcal O_{L_w})$. The source fixes $p$ odd and unramified
in $L$, making this principal congruence subgroup torsion-free.
For example, if $h=1+p^aC\ne1$ with some entry of $C$ a unit, a
prime-to-$p$ power has the same leading valuation, while
$h^p=1+p^{a+1}C'$ with $C'\equiv C\pmod p$.
Thus the stabilizer is trivial. The place $v_0$ is unchanged by
every allowed auxiliary level.

Put $A_a=A/\mathfrak m^a$, $a\geq1$, and choose $c$ such that
$\theta(T_c)\equiv1\pmod{\mathfrak m^a}$. The space $E_{a,c}$ of
functions with the original transformation laws and exact
$\overline N(c)$ invariance is finite free over $A_a$: evaluation
on the finite double-coset representatives is an isomorphism.
Gaussian elimination gives
$I(c)=NT\overline N(c)=\overline N(c)TN$, and its upper/lower
interchange changes the torus component only by $T_c$. This
both proves the stated transformation rule over $A_a$ and
constructs integral lifts of arbitrary representative values:
use $\tau(u)^{-1}\theta(t(b))^{-1}$, with $t(b)$ the torus
component in the latter factorization. Right multiplication by
$\overline N(d)$ changes $t(b)$ by $T_d$, so these lifts have
the required uniform continuity.

The original double-coset formula for $U$ lowers
$\overline N(d+1)$ invariance to $\overline N(d)$ invariance
when $d\geq c$ and the vector is fixed by $T_c$. It follows
from $z^{-1}\overline N(d)z\subset\overline N(d+1)$ and the
same upper/lower interchange. Consequently $U$ preserves
$E_{a,c}$ and
$U^{c'-c}E_{a,c'}\subset E_{a,c}$ for $c'\geq c$.
The invertible Fitting summand $M_a$ of $U$ on $E_{a,c}$ is
therefore independent of $c$. It is a finite free $A_a$-module,
and its projector is a sufficiently large factorial power of $U$.
Choosing the same $c$ at two precisions gives, in the same
double-coset bases,
\[
 M_a\otimes_{A_a}A_b=M_b\qquad(1\leq b\leq a).
 \tag{OD3}
\]
Thus all ranks agree.

For each $f\in S(K,A)$ and each $a$, a sufficiently high
$U$-iterate lies in $E_{a,c}$ modulo $\mathfrak m^a$.
Hence $U^{j!}f$ has compatible stabilized reductions, defining
an $A$-linear projector $e$ on $S(K,A)$. Completeness preserves
the covariance laws; continuity modulo each $\mathfrak m^a$
preserves uniform continuity. The fixed space of $e$ is exactly
(OD1). The integral sections give surjectivity onto each $M_a$,
and a compatible sequence in the $M_a$ defines an original
continuous automorphic function. Therefore
$M(K,A)=\varprojlim_aM_a$. Compatible bases in (OD3) prove
finite freeness and the asserted filtration equality.
The construction commutes with tame translations and commuting
Hecke operators.

Now put $M=M(K',A)$. Neatness makes the $\Delta$-action on
the original double cosets at level $K'I(0)$ free. Choose one
representative per orbit, let $\chi$ be its characteristic
function, and define $\phi=e\circ M_\chi$ on $M$.
The identity $\sum_\delta\chi(g\delta)=1$ gives
\[
 \sum_{\delta\in\Delta}\delta\phi\delta^{-1}=1_M.
 \tag{OD4}
\]
It follows that
$m\mapsto\sum_\delta\delta\otimes\phi(\delta^{-1}m)$
splits the action map $A[\Delta]\otimes_AM\to M$
$A[\Delta]$-linearly. The domain is finite free, proving
projectivity. For invariant $m$, (OD4) gives
$\operatorname{Tr}(\phi m)=m$. If $\operatorname{Tr}(m)=0$,
it gives
$m=\sum_\delta(\delta-1)\phi(\delta^{-1}m)$.
These are precisely (OD2), with no division by $|\Delta|$.

The commutative Hecke algebra in $\operatorname{End}_A(M)$
is finite over the complete Noetherian ring $A$. Its local
factors supply the idempotent $e_{\mathfrak h}$ for the
prescribed residual eigenvalues, or zero if they are absent.
This idempotent commutes with $\Delta$ and restricts to the
corresponding factor on the invariant space. Replacing $\phi$
by $e_{\mathfrak h}\phi$ on $e_{\mathfrak h}M$ preserves
(OD4), proving the localized assertion
\cite{alpo-ordinary-diamond}.
\end{proof}

At the actual auxiliary places $w_j$, the prescribed quotient is
\[
 \Delta=\prod_j(\mathbb Z/\ell^{a_j}\mathbb Z)^n,\qquad
 a_j=v_\ell(\mathrm Nw_j-1),
\]
obtained from the diagonal residue map by killing the
prime-to-$\ell$ subgroup. The fixed tame subgroup has identity
projection at each auxiliary place, while the original local
sections project to their specified coordinates. Thus this is
also the actual tame quotient, with no absorbed diamond factor.
With the original generators
$\delta_{ji}$, its group ring is exactly
\[
 A[\Delta]\simeq
 A[[Y_{ji}]]/
   \bigl((1+Y_{ji})^{\ell^{a_j}}-1\bigr)_{j,i},
 \qquad Y_{ji}\longmapsto\delta_{ji}-1.
 \tag{OD5}
\]
Division by the displayed monic distinguished polynomials proves
the isomorphism. Modulo $\mathfrak m$ this is a truncated polynomial
ring with its unique augmentation maximal ideal. Thus (OD5) is
local, and the modules in the proposition are free over it.
Any later commuting spectral idempotent gives a direct summand,
again free over the same ring; its augmentation quotient is the
corresponding projected Iwahori module by (OD2).
These are the original module and trace clauses of the
Taylor--Wiles construction. The smooth union over all auxiliary
levels has exactly these finite invariant modules. Here the
auxiliary set, coefficient representation, character and residual
eigenvalues are fixed. The Hecke family is the original spherical
family outside that fixed set together with the positive diagonal
operators above $\ell$; it commutes with auxiliary translations
at every level. No new Hecke condition is introduced in forming
the union.

The statement is valid for every $a_j\geq0$, including the
original requirement $a_j\geq N$, without a new exponent loss.
It neither bounds ranks uniformly as the arithmetic datum moves
nor proves that an arbitrary residual component is nonzero.
The original universal-determinant action, separating classical
specializations with local Galois compatibility, and the
spherical/Iwahori comparison maps remain to be supplied.
